\documentclass[11pt]{article}

% Packages for XeLaTeX
\usepackage{fontspec}
\usepackage{amsmath,amssymb,amsthm}
\usepackage{geometry}
\usepackage{hyperref}
\usepackage{enumitem}
\usepackage{mathtools}
\usepackage{bm}
\usepackage{framed}
\usepackage{xcolor}
\usepackage{algorithm}
\usepackage{algorithmic}
\usepackage{tikz}
\usetikzlibrary{calc}
\usepackage{eso-pic}
\usepackage[useregional]{datetime2}

\geometry{margin=1in}

% Running margin label
\newcommand{\mymarginlabel}{David Ewing dew59@uclive.ac.nz | \DTMnow}
\AddToShipoutPictureBG{%
  \begin{tikzpicture}[remember picture,overlay]
    \node[rotate=90, anchor=north]
      at ($(current page.north west)+(1.4cm,-13cm)$)
      {\scriptsize\ttfamily \mymarginlabel};
  \end{tikzpicture}%
}

% Custom commands
\newcommand{\E}{\mathbb{E}}
\newcommand{\R}{\mathbb{R}}
\newcommand{\N}{\mathcal{N}}
\newcommand{\ELBO}{\text{ELBO}}
\newcommand{\tr}{\text{tr}}
\setlength{\parindent}{0pt}

\title{ 
Optimisation Methods to Address\\
Variational Bayesian\\
Posterior Variance Collapse
}
\author{
  David Ewing (82171165)\\
%ENEL445  \\
%University of Canterbury
}
\date{3 May 2026}

\begin{document}

\maketitle

\begin{abstract}
\setlength{\parindent}{0pt}
\setlength{\parskip}{0.7em}

Variational Inference (VI) recasts Bayesian posterior computation as an optimisation problem, making it a natural testbed for comparing the
gradient-based methods taught in ENEL445. This document provides the mathematical derivations for VI applied to Bayesian linear regression under
a mean-field Normal--Gamma approximation.

We derive the Evidence Lower Bound (ELBO) in closed form and present the
analytical Coordinate Ascent VI (CAVI) solution as a baseline. For
gradient-based methods---gradient ascent, Newton's method, and BFGS---we
derive full gradient and Hessian expressions for all variational parameters
($\bm{\mu}_\beta$, $\bm{\Sigma}_\beta$, $a_e$, $b_e$), and formulate a
quadratic penalty method to enforce minimum-variance constraints that
prevent posterior variance collapse under the zero-forcing KL divergence.

The entry conditions are stated explicitly for each method.   For line-search methods, a further
condition must hold before each step is accepted: the search direction
$\bm{d}^{(t)}$ must satisfy $\nabla\ELBO(\bm{\phi}^{(t)})^T \bm{d}^{(t)}
> 0$ (ascent direction), ensuring the line search is well-posed.

The step size $\alpha^{(t)}$ is determined by backtracking line
search. Gradient ascent uses the Armijo sufficient-increase condition. BFGS uses the stronger Wolfe conditions---sufficient
increase plus a curvature condition---to guarantee the gradient-difference
vector $\bm{y}_t^T \bm{s}_t > 0$, which is required for the BFGS update
to maintain positive definiteness of the approximate inverse Hessian.

Exit conditions are defined uniformly across all methods.
Iteration terminates when any of the following hold: (i) gradient norm
$\lVert\nabla\ELBO\rVert < \varepsilon_g$; (ii) relative ELBO change
$\lvert\Delta\ELBO\rvert / \lvert\ELBO\rvert < \varepsilon_r$; (iii)
parameter change $\lVert\Delta\bm{\phi}\rVert < \varepsilon_\phi$; or
(iv) a maximum iteration count is reached. CAVI additionally monitors
$\max_i \lvert\phi_i^{(t+1)} - \phi_i^{(t)}\rvert /
\lvert\phi_i^{(t+1)}\rvert < \varepsilon$.

Convergence rates and per-iteration costs are compared across all methods.
BFGS offers the best practical balance of convergence speed and
computational cost; CAVI remains competitive when conjugate structure can
be exploited analytically.
\end{abstract}
\newpage
\section*{Introduction}

Bayesian inference provides principled uncertainty quantification but requires computing intractable posterior distributions. For most real-world models, the posterior integral,

\begin{equation*}
    p(\bm{\theta} \mid \bm{y}) = \frac{p(\bm{y} \mid \bm{\theta})\,p(\bm{\theta})}{\displaystyle\int p(\bm{y} \mid \bm{\theta})\,p(\bm{\theta})\, d\bm{\theta}},
\end{equation*}
\noindent 
cannot be evaluated analytically. 
Variational Inference (VI) converts this inference problem into an optimisation problem, making it directly amenable to the gradient-based methods in this course. \vspace{2em}

Instead of computing the true posterior, VI finds the ``best'' approximation $q(\bm{\theta})$ from a tractable family $Q(\bm{\theta})$
by maximising the Evidence Lower Bound (ELBO). 
This transforms intractable integration into tractable optimisation---a natural testbed for the gradient-based methods. \vspace{2em} 

A well-known failure mode of mean-field VI is \emph{Posterior Variance Collapse}: 
the optimised approximation $q(\bm{\theta})$ is systematically overconfident, 
underestimating posterior uncertainty. This is a direct consequence of the KL direction minimised during ELBO maximisation, 
compounded by the mean-field factorisation assumption. \footnote{See \ref{app:vb-collapse} for the full mechanistic explanation.}  \vspace{2em} 

This paper addresses variance collapse via two complementary strategies, 
\begin{itemize}
\item The first applies new optimisation methods to the \emph{standard} ELBO objective: gradient ascent, Newton's method, and BFGS are compared to determine whether the choice of optimiser affects dispersion quality.
\item The second modifies the objective function itself, augmenting the ELBO with an entropy regularisation term $\lambda H[q(\bm{\beta})]$ to create a new objective $\mathcal{L}_{\mathrm{reg}}(\bm{\phi}) = \mathcal{L}(\bm{\phi}) + \lambda H[q(\bm{\beta})]$ that explicitly incentivises broader posteriors.
\end{itemize}
Both the optimiser-comparison and modified-ELBO strategies are explored.  \vspace{2em} 

Throughout all gradient-based methods, the variational covariance $\bm{\Sigma}_\beta$ must remain symmetric positive definite at every iteration.\footnote{This structural constraint is satisfied automatically via Cholesky reparameterisation; see Appendix~A. For the role of symmetric positive definite covariance matrices in exponential families more broadly, see Wainwright \& Jordan (2008), \S3.4, DOI: \href{https://doi.org/10.1561/2200000001}{10.1561/2200000001}.} Standard gradient steps in Euclidean space can violate this requirement, leading to numerical failure.
 The solution is to parameterise the optimisation over the Cholesky factor $\bm{L}$ rather than $\bm{\Sigma}_\beta$ directly.
 Cholesky reparameterisation is a numerical implementation technique .   Cholesky is the domain-specific choice for satisfying the constraint $\bm{\Sigma}_\beta \succ 0$ when applying those methods to this problem. The set of symmetric positive definite matrices is not a flat subspace of $\R^{p \times p}$: it has a curved boundary, and a straight-line gradient step can cross outside it, producing an invalid covariance matrix. Cholesky reparameterisation handles this by converting the constraint into a smooth bijection between the valid set and unconstrained $\R^{p(p+1)/2}$, so standard gradient steps always remain valid. It is standard practice in variational inference and Bayesian computation.


\section*{Key Notation}

\begin{center}
\begin{tabular}{ll}
\hline
Symbol & Meaning \\
\hline
$\bm{\beta}$ & Regression coefficient vector \\
$\bm{u}$ & Random effects vector (hierarchical model) \\
$\bm{\varepsilon}$ & Observation noise vector \\
$\tau_e$ & Observation-noise precision \\
$\tau_u$ & Random-effects precision \\
$\bm{\mu}_\beta$ & Mean of variational posterior $q(\bm{\beta})$ \\
$\bm{\Sigma}_\beta$ & Covariance of variational posterior $q(\bm{\beta})$ \\
$\bm{L}$ & Lower-triangular Cholesky factor of $\bm{\Sigma}_\beta$ \\
$\eta_a, \eta_b$ & Log-space parameters for $a_e, b_e$ \\
\hline
\end{tabular}
\end{center}
\newpage
\section*{Optimisation Problem Formulation}

Having established that VI converts posterior inference into optimisation, the problem must be stated precisely before any method can be applied. This section identifies the decision variables, the objective function, and the constraints, following the standard form used throughout the course. A well-posed formulation is necessary both for implementation and for meaningful comparison across methods: gradient ascent, Newton's method, and BFGS all operate on the same objective and respect the same constraints, so differences in performance reflect the methods themselves rather than problem differences.

\subsection*{Decision Variables}

For Bayesian linear regression with $p$ predictors, the variational parameters (decision variables) are:
\begin{itemize}
\item $\bm{\mu}_\beta \in \R^p$: posterior mean of regression coefficients
\item $\bm{\Sigma}_\beta \in \R^{p \times p}$: posterior covariance (symmetric positive definite\footnote{A matrix $\bm{A}$ is \textbf{positive definite} 
if $\bm{v}^T\bm{A}\bm{v} > 0$ for every non-zero vector $\bm{v}$. Intuitively: the matrix amplifies every 
direction --- there is no direction in which it squashes things to zero or flips them negative. For a
 covariance matrix this means every direction in parameter space has genuine positive uncertainty. 
 Written $\bm{A} \succ 0$.})
\item $a_e, b_e > 0$: shape and rate parameters of the variational Gamma factor for precision
\end{itemize}

Here, $a_e$ and $b_e$ parameterise $q(\tau_e)=\mathrm{Gamma}(a_e,b_e)$; the model prior on $\tau_e$ is specified in the Test Case section.

Total dimension: $d = p + \frac{p(p+1)}{2} + 2 = 11$ variables for $p=3$ predictors.

\subsection*{Objective Function}

Maximise the Evidence Lower Bound (ELBO):

\begin{align}
\ELBO(\bm{\phi})
&= \E_q[\log p(\bm{y}, \bm{\beta}, \tau_e)] + H(q) \\
&= \E_q[\log p(\bm{y}, \bm{\beta}, \tau_e)] - \E_q[\log q(\bm{\beta}, \tau_e)] .
\end{align}

where $\bm{\phi} = (\bm{\mu}_\beta, \bm{\Sigma}_\beta, a_e, b_e)$ are the variational parameters.

\subsection*{Constraints}

\begin{itemize}
\item $\bm{\Sigma}_\beta \succ 0$: Covariance matrix must be positive definite (see also: positive semidefinite, $\bm{A} \succeq 0$\footnote{\textbf{Positive semidefinite} ($\bm{A} \succeq 0$) allows $\bm{v}^T\bm{A}\bm{v} = 0$ for some non-zero $\bm{v}$, meaning some combination of coefficients has \emph{exactly} zero uncertainty --- the distribution is flat in that direction. A valid covariance matrix only needs to be positive semidefinite, but a semidefinite (not strictly definite) matrix is non-invertible\footnote{A matrix is \textbf{invertible} if $\bm{A}^{-1}$ exists, meaning the equation $\bm{A}\bm{x}=\bm{b}$ has a unique solution for any $\bm{b}$. Equivalently: none of the eigenvalues are zero. The precision matrix $\bm{\Sigma}_\beta^{-1}$ appears in the ELBO and CAVI updates; if $\bm{\Sigma}_\beta$ is not invertible these computations fail.} and leads to degenerate directions.\footnote{A \textbf{degenerate direction} is a direction $\bm{v}$ in which the distribution has zero spread --- the posterior is completely certain that $\bm{v}^T\bm{\beta}$ takes one exact value with no uncertainty. Geometrically: the probability ellipsoid is flat (squashed to zero thickness) along $\bm{v}$. This corresponds to a zero eigenvalue and makes the covariance matrix singular (non-invertible). It usually signals numerical collapse or that the data perfectly determines some linear combination of $\bm{\beta}$.}})
\item $a_e, b_e > 0$: Gamma parameters must be positive
\item Optional: Minimum variance constraints $\sigma^2_{\beta_j} \geq \sigma^2_{\text{min}}$ to prevent underdispersion
\end{itemize}

\textbf{Constraint Handling via Reparameterisation:}

To enforce $\bm{\Sigma}_\beta \succ 0$, use Cholesky decomposition:
\begin{align}
\bm{\Sigma}_\beta &= \bm{L}\bm{L}^T
\end{align}
where $\bm{L} \in \R^{p \times p}$ is lower triangular with positive diagonal entries. The unconstrained parameters are:
\begin{itemize}
\item $L_{ii} = e^{\ell_{ii}}$ for $i=1,\ldots,p$ (exponentiate to ensure positivity)
\item $L_{ij} = \ell_{ij}$ for $i > j$ (off-diagonal entries unconstrained)
\end{itemize}

Total Cholesky parameters: $\frac{p(p+1)}{2}$ elements of $\bm{\ell}$.

To enforce $a_e, b_e > 0$, use log-space parameterisation:
\begin{align}
a_e &= e^{\eta_a}, \quad b_e = e^{\eta_b}
\end{align}
where $\eta_a, \eta_b \in \R$ are unconstrained.

\textbf{Unconstrained parameter vector:} $\bm{\xi} = (\bm{\mu}_\beta, \bm{\ell}, \eta_a, \eta_b) \in \R^d$ with $d = p + \frac{p(p+1)}{2} + 2 = 11$ for $p=3$.

\subsection*{Computational Characteristics}

\begin{itemize}
\item Gradient computation: $O(np^2)$ per evaluation (dominated by matrix operations)
\item Hessian computation: $O(np^2 + d^3)$ for Newton's method
\item Parameters exhibit coupling (changing $\bm{\mu}_\beta$ affects optimal $\bm{\Sigma}_\beta$)
\item Must handle special functions: digamma $\psi(\cdot)$, trigamma $\psi'(\cdot)$ for Gamma distribution
\end{itemize}
\newpage
\section*{Test Case: Bayesian Linear Regression}

The test case uses Bayesian linear regression with:
\begin{align}
\bm{y} &\sim \N(\bm{X}\bm{\beta}, \tau_e^{-1} \bm{I}_n) \quad \text{(likelihood)}\\
\bm{\beta} &\sim \N(\bm{0}, \lambda_\beta^{-1} \bm{I}_p) \quad \text{(prior on coefficients)}\\
\tau_e &\sim \text{Gamma}(\alpha_e, \gamma_e) \quad \text{(prior on precision)}
\end{align}

Mean-field approximation:
\begin{align}
q(\bm{\beta}, \tau_e) &= q(\bm{\beta}) \, q(\tau_e)\\
q(\bm{\beta}) &= \N(\bm{\mu}_\beta, \bm{\Sigma}_\beta)\\
q(\tau_e) &= \text{Gamma}(a_e, b_e)
\end{align}

Variational parameters: $\bm{\phi} = (\bm{\mu}_\beta, \bm{\Sigma}_\beta, a_e, b_e)$ with dimension $d = p + \frac{p(p+1)}{2} + 2$

For $p=3$ predictors: $d = 3 + 6 + 2 = 11$ parameters to optimise.
 
\section*{Overview of Optimisation Methods}

\subsection*{Methods considered}

\begin{itemize}
\item CAVI: Domain-specific baseline exploiting problem structure
\item Gradient Ascent: First-order baseline method using gradient information and line search
\item Newton's Method: Second-order method with quadratic local convergence near a well-conditioned optimum
\item BFGS: Quasi-Newton method that approximates Hessian information; practical general-purpose method for smooth unconstrained problems
\end{itemize}

\subsection*{Baseline Method: CAVI\footnote{Blei, Kucukelbir, and McAuliffe (2017), Variational Inference: A Review for Statisticians (Journal of the American Statistical Association).}}

Coordinate Ascent Variational Inference (CAVI) is VI-specific. It updates each variational parameter block sequentially using closed-form solutions:

\begin{align}
\bm{\Sigma}_\beta &\leftarrow \left(\frac{a_e}{b_e} \bm{X}^T\bm{X} + \lambda_\beta \bm{I}_p\right)^{-1}\\
\bm{\mu}_\beta &\leftarrow \frac{a_e}{b_e} \bm{\Sigma}_\beta \bm{X}^T\bm{y}\\
a_e &\leftarrow \alpha_e + \frac{n}{2}\\
b_e &\leftarrow \gamma_e + \frac{1}{2}\left[\|\bm{y} - \bm{X}\bm{\mu}_\beta\|^2 + \tr(\bm{X}^T\bm{X}\bm{\Sigma}_\beta)\right]
\end{align}

Properties: Monotonic ELBO increase\footnote{Monotonic ELBO increase means the objective is non-decreasing across iterations: $\mathcal{L}^{(t+1)} \ge \mathcal{L}^{(t)}$. In CAVI, each coordinate update maximises the ELBO with respect to one variational factor while holding the others fixed, so each step cannot decrease the ELBO.}, linear convergence, $O(np^2)$ per iteration. Provides analytical benchmark for validation.

\newpage 

\subsection*{Mapping ENEL445 Methods to VI Problem}

\begin{enumerate}[leftmargin=*]
\item Parameter Ascent (Course: gradient descent)
\begin{itemize}
\item Standard iterative update: $\bm{\phi}^{(t+1)} = \bm{\phi}^{(t)} + \alpha_t \nabla \ELBO(\bm{\phi}^{(t)})$
\item Line search from Chapter 3.2 (bisection/Wolfe conditions)
\item Stopping criteria: $\|\nabla \ELBO\| < \epsilon$
\item Convergence characteristics: linear rate expected for strongly concave ELBO
\end{itemize}

\item Newton's Method   
\begin{itemize}
\item Standard Newton update: $\bm{\phi}^{(t+1)} = \bm{\phi}^{(t)} + \alpha_t [\nabla^2 \ELBO]^{-1} \nabla \ELBO$
\item Requires computing and inverting the Hessian matrix
\item Computational cost: $O(d^3)$ per iteration for Hessian inversion
\item Convergence characteristics: quadratic convergence near optimum (when Hessian is negative definite)
\item Challenge: Hessian may not be negative definite away from the optimum
\end{itemize}

\item BFGS Quasi-Newton (one of five required methods from Project 2)
\begin{itemize}
\item BFGS update formula builds inverse Hessian approximation
\item Uses gradient differences: $\bm{s}_k = \bm{\phi}^{(k)} - \bm{\phi}^{(k-1)}$, $\bm{y}_k = \nabla \ELBO(\bm{\phi}^{(k)}) - \nabla \ELBO(\bm{\phi}^{(k-1)})$ 
\item Computational cost: $O(d^2)$ per iteration
\item Convergence characteristics: superlinear convergence expected
\end{itemize}
\end{enumerate}

Comparison will reveal whether general optimisation methods can compete with or outperform problem-specific algorithms.

\begin{itemize}
\item[\checkmark] Complete ELBO derivation for Bayesian linear regression
\item[\checkmark] CAVI closed-form update equations with matrix inversion lemma
\item[\checkmark] Gradient expressions: $\nabla_{\bm{\mu}_\beta} \ELBO$, $\nabla_{\bm{\Sigma}_\beta} \ELBO$, $\nabla_{a_e} \ELBO$, $\nabla_{b_e} \ELBO$
\item[\checkmark] Gradient transformation for Cholesky parameterisation: $\nabla_{\bm{\ell}} \ELBO$
\item[\checkmark] Gradient transformation for log-space parameterisation: $\nabla_{\eta_a} \ELBO$, $\nabla_{\eta_b} \ELBO$
\item[\checkmark] Hessian block structure: diagonal and mixed blocks
\item[\checkmark] BFGS update formulation with Woodbury identity
\item[\checkmark] Initialisation strategy with principled defaults
\item[\checkmark] Line search algorithm specifications (Armijo backtracking)
\item[\checkmark] Newton regularisation strategy for indefinite Hessian
\item[\checkmark] Convergence criteria with multiple stopping conditions
\end{itemize}

All derivations documented in \texttt{vi\_optimisation\_solution.pdf} with complete algebraic steps.

(Truncated for brevity in file creation)
